% GENERATED FILE. Edit canonical lessons, then run npm run generate:books.
% canonical-source-hash: fnv1a64:501ed455dfa3ef20
% canonical-lessons: GU-C464-sanu, GU-C464-vivek, GU-C464-dahyu, GU-C464-thavku, GU-C464-gathvu

\chapter{A Sensible Head}
\label{ch:gu-a-sensible-head}
\hlchaptermodality{464}

\begin{chapteropening}
\textbf{By the end of this chapter:} \emph{I can describe a wise and sensible person.}

Complete the last lesson of chapter 464: I can describe a wise and sensible person.
\end{chapteropening}

\section[\texorpdfstring{\gu{શાણું} (śāṇũ)}{શાણું (śāṇũ)}]{\gu{શાણું} (śāṇũ) — wise; discreet}

\label{lesson:GU-C464-sanu}



\begin{quote}
Before the new one: say the Gujarati for to babble, then the Gujarati for cheerful.
\end{quote}

\subsection*{You'll want to know: \gu{શાણું}}
\textbf{\gu{શાણું}} — \emph{śāṇũ} — \textquotedblleft{}wise; discreet\textquotedblright{}.

\begin{quote}
\gu{શાણો} \gu{માણસ} \gu{ગુસ્સામાં} \gu{કંઈ} \gu{બોલતો} \gu{નથી}.
\end{quote}

\begin{quote}
\emph{śāṇo māṇas gussāmā̃ kãī bolto nathī.}
\end{quote}

\begin{quote}
A wise man says nothing in anger.
\end{quote}

\begin{cousinweb}
\textbf{\gu{શાણું}} (\emph{śāṇũ}) goes back to a reconstructed \emph{*sajāna-}, \textquotedblleft{}clever\textquotedblright{}; CDIAL lists the Gujarati word, \textquotedblleft{}wise, discreet, courteous\textquotedblright{} (CDIAL 13088).
\end{cousinweb}

\begin{grammarlens}[title={What you've built}]
One more word for describing what you lived and felt.
\end{grammarlens}

\subsection*{Your turn}
\begin{itemize}
\raggedright
  \item \emph{Say it:} \emph{śāṇũ}
  \item \emph{Say it:} \emph{śāṇũ}, once more
  \item \emph{Say it:} \emph{śāṇo māṇas gussāmā̃ kãī bolto nathī.} — A wise man says nothing in anger.
  \item \emph{Recall:} say the Gujarati for to babble, then the Gujarati for cheerful, then say \emph{śāṇũ} again
\end{itemize}

\begin{tcolorbox}[breakable,colback=teal!4,colframe=teal!35!black,title={Before you move on}]
What does \gu{શાણું} mean? (\textquotedblleft{}Wise; discreet\textquotedblright{}.) Say it once more.
\end{tcolorbox}
\section[\texorpdfstring{\gu{વિવેક} (vivek)}{વિવેક (vivek)}]{\gu{વિવેક} (vivek) — good sense; courtesy}

\label{lesson:GU-C464-vivek}



\begin{quote}
Before the new one: say the Gujarati for cheerful, then the Gujarati for discreet.
\end{quote}

\subsection*{You'll want to know: \gu{વિવેક}}
\textbf{\gu{વિવેક}} — \emph{vivek} — \textquotedblleft{}good sense; courtesy\textquotedblright{}.

\begin{quote}
\gu{મહેમાન} \gu{સાથે} \gu{વિવેકથી} \gu{બોલો}.
\end{quote}

\begin{quote}
\emph{mahemān sāthe vivekthī bolo.}
\end{quote}

\begin{quote}
Speak to guests with courtesy.
\end{quote}

A masculine noun.

\begin{cousinweb}
\textbf{\gu{વિવેક}} (\emph{vivek}) is a tatsama, Sanskrit \emph{viveka}, \textquotedblleft{}discrimination, consideration\textquotedblright{}, as Monier-Williams records it; Gujarati uses it for good sense and good manners.
\end{cousinweb}

\begin{grammarlens}[title={What you've built}]
One more word for describing what you lived and felt.
\end{grammarlens}

\subsection*{Your turn}
\begin{itemize}
\raggedright
  \item \emph{Say it:} \emph{vivek}
  \item \emph{Say it:} \emph{vivek}, once more
  \item \emph{Say it:} \emph{mahemān sāthe vivekthī bolo.} — Speak to guests with courtesy.
  \item \emph{Recall:} say the Gujarati for cheerful, then the Gujarati for discreet, then say \emph{vivek} again
\end{itemize}

\begin{tcolorbox}[breakable,colback=teal!4,colframe=teal!35!black,title={Before you move on}]
What does \gu{વિવેક} mean? (\textquotedblleft{}Good sense; courtesy\textquotedblright{}.) Say it once more.
\end{tcolorbox}
\section[\texorpdfstring{\gu{ડાહ્યું} (ḍāhyũ)}{ડાહ્યું (ḍāhyũ)}]{\gu{ડાહ્યું} (ḍāhyũ) — well-behaved; sensible}

\label{lesson:GU-C464-dahyu}



\begin{quote}
Before the new one: say the Gujarati for discreet, then the Gujarati for good sense.
\end{quote}

\subsection*{You'll want to know: \gu{ડાહ્યું}}
\textbf{\gu{ડાહ્યું}} — \emph{ḍāhyũ} — \textquotedblleft{}well-behaved; sensible\textquotedblright{}.

\begin{quote}
\gu{ડાહ્યું} \gu{બાળક} \gu{મહેમાનો} \gu{સામે} \gu{શાંત} \gu{બેઠું}.
\end{quote}

\begin{quote}
\emph{ḍāhyũ bāḷak mahemāno sāme śānt beṭhũ.}
\end{quote}

\begin{quote}
The well-behaved child sat quietly in front of the guests.
\end{quote}

\begin{cousinweb}
\textbf{\gu{ડાહ્યું}} (\emph{ḍāhyũ}) goes back to Sanskrit \emph{dāsa}, \textquotedblleft{}a knowing man\textquotedblright{}, a word of the old word lists; CDIAL lists the Gujarati word, \textquotedblleft{}wise\textquotedblright{} (CDIAL 6317).
\end{cousinweb}

\begin{grammarlens}[title={What you've built}]
One more word for describing what you lived and felt.
\end{grammarlens}

\subsection*{Your turn}
\begin{itemize}
\raggedright
  \item \emph{Say it:} \emph{ḍāhyũ}
  \item \emph{Say it:} \emph{ḍāhyũ}, once more
  \item \emph{Say it:} \emph{ḍāhyũ bāḷak mahemāno sāme śānt beṭhũ.} — The well-behaved child sat quietly in front of the guests.
  \item \emph{Recall:} say the Gujarati for discreet, then the Gujarati for good sense, then say \emph{ḍāhyũ} again
\end{itemize}

\begin{tcolorbox}[breakable,colback=teal!4,colframe=teal!35!black,title={Before you move on}]
What does \gu{ડાહ્યું} mean? (\textquotedblleft{}Well-behaved; sensible\textquotedblright{}.) Say it once more.
\end{tcolorbox}
\section[\texorpdfstring{\gu{ઠાવકું} (ṭhāvkũ)}{ઠાવકું (ṭhāvkũ)}]{\gu{ઠાવકું} (ṭhāvkũ) — composed; sober}

\label{lesson:GU-C464-thavku}



\begin{quote}
Before the new one: say the Gujarati for good sense, then the Gujarati for well-behaved.
\end{quote}

\subsection*{You'll want to know: \gu{ઠાવકું}}
\textbf{\gu{ઠાવકું}} — \emph{ṭhāvkũ} — \textquotedblleft{}composed; sober\textquotedblright{}.

\begin{quote}
\gu{આટલી} \gu{નાની} \gu{ઉંમરે} \gu{પણ} \gu{મીરા} \gu{ઘણી} \gu{ઠાવકી} \gu{છે}.
\end{quote}

\begin{quote}
\emph{āṭlī nānī ummre paṇ Mīrā ghaṇī ṭhāvkī che.}
\end{quote}

\begin{quote}
Even at so young an age Mira is very composed.
\end{quote}

\begin{cousinweb}
\textbf{\gu{ઠાવકું}} (\emph{ṭhāvkũ}) is inherited: CDIAL lists the Gujarati word, \textquotedblleft{}steady, firm, discreet\textquotedblright{}, under Sanskrit \emph{sthātavya}, \textquotedblleft{}to be stood in\textquotedblright{} (CDIAL 13751): someone who stands firm.
\end{cousinweb}

\begin{grammarlens}[title={What you've built}]
One more word for describing what you lived and felt.
\end{grammarlens}

\subsection*{Your turn}
\begin{itemize}
\raggedright
  \item \emph{Say it:} \emph{ṭhāvkũ}
  \item \emph{Say it:} \emph{ṭhāvkũ}, once more
  \item \emph{Say it:} \emph{āṭlī nānī ummre paṇ Mīrā ghaṇī ṭhāvkī che.} — Even at so young an age Mira is very composed.
  \item \emph{Recall:} say the Gujarati for good sense, then the Gujarati for well-behaved, then say \emph{ṭhāvkũ} again
\end{itemize}

\begin{tcolorbox}[breakable,colback=teal!4,colframe=teal!35!black,title={Before you move on}]
What does \gu{ઠાવકું} mean? (\textquotedblleft{}Composed; sober\textquotedblright{}.) Say it once more.
\end{tcolorbox}
\section[\texorpdfstring{\gu{ગાંઠવું} (gā̃ṭhvũ)}{ગાંઠવું (gā̃ṭhvũ)}]{\gu{ગાંઠવું} (gā̃ṭhvũ) — to heed; to pay attention to (someone)}

\label{lesson:GU-C464-gathvu}



\begin{quote}
Before the new one: say the Gujarati for well-behaved, then the Gujarati for composed.
\end{quote}

\subsection*{You'll want to know: \gu{ગાંઠવું}}
\textbf{\gu{ગાંઠવું}} — \emph{gā̃ṭhvũ} — \textquotedblleft{}to heed; to pay attention to (someone)\textquotedblright{}.

\begin{quote}
\gu{તોફાની} \gu{દીકરો} \gu{કોઈનું} \gu{ગાંઠતો} \gu{નથી}.
\end{quote}

\begin{quote}
\emph{tophānī dīkro koīnũ gā̃ṭhto nathī.}
\end{quote}

\begin{quote}
The naughty son heeds no one.
\end{quote}

Past, with \gu{તેણે} (he or she …): \gu{તેણે} \gu{કોઈનું} \gu{ન} \gu{ગાંઠ્યું}, \emph{teṇe koīnũ na gā̃ṭhyũ}.

\begin{cousinweb}
\textbf{\gu{ગાંઠવું}} (\emph{gā̃ṭhvũ}) is inherited from Sanskrit \emph{granthayati}, \textquotedblleft{}ties\textquotedblright{}; CDIAL lists the Gujarati verb, \textquotedblleft{}to tie\textquotedblright{} (CDIAL 4353), and Gujarati also uses it, mostly with a negative, for heeding someone.
\end{cousinweb}

\begin{grammarlens}[title={What you've built}]
One more word for describing what you lived and felt.
\end{grammarlens}

\subsection*{Your turn}
\begin{itemize}
\raggedright
  \item \emph{Say it:} \emph{gā̃ṭhvũ}
  \item \emph{Say it:} \emph{gā̃ṭhvũ}, once more
  \item \emph{Say it:} \emph{tophānī dīkro koīnũ gā̃ṭhto nathī.} — The naughty son heeds no one.
  \item \emph{Recall:} say the Gujarati for well-behaved, then the Gujarati for composed, then say \emph{gā̃ṭhvũ} again
\end{itemize}

\begin{tcolorbox}[breakable,colback=teal!4,colframe=teal!35!black,title={Before you move on}]
What does \gu{ગાંઠવું} mean? (\textquotedblleft{}To heed; to pay attention to (someone)\textquotedblright{}.) Say it once more.
\end{tcolorbox}
